
\begin{table}[t]\centering\footnotesize
\caption{Comparison with the closest spline-collocation approaches, to the best of our knowledge.
``Spectrum'' refers to the one-dimensional collocation operator used by fast diagonalization or matrix
decomposition.}\label{tab:closest}
\begin{tabular}{p{0.14\textwidth}p{0.18\textwidth}p{0.13\textwidth}p{0.235\textwidth}p{0.175\textwidth}}
\toprule
Method & Trial space and collocation points & Grids & Spectrum & Fast solver\\
\midrule
Classical HWCM~\cite{lepik2011solving,oruc2019nonuniform} & integrated Haar ($C^1$ quadratics); midpoints & uniform or a~priori graded & not analysed & none (dense $N^2\times N^2$ solve)\\
Hermite-cubic OSC~\cite{dyksen1987tensor,sun1999spectral,bialecki1995matrix} & $C^1$ cubics; two Gauss points per cell & partitions as in the cited works & real, distinct eigenvalues established & matrix decomposition; FFT on uniform grids\\
QSC~\cite{houstis1988quadratic,christara2002fast,christara2006optimal} & $C^1$ quadratics; midpoints (optimal variants perturbed) & uniform; non-uniform analysed for convergence & symmetric, trigonometric eigenvectors on uniform grids & FFT and matrix decomposition on uniform grids\\
Quadratic IGA collocation~\cite{auricchio2010isogeometric,schillinger2013isogeometric,delorenzis2015isogeometric} & B-splines; Greville or superconvergent points & arbitrary knots; local refinement & not established for collocation; Galerkin pencils symmetric~\cite{sangalli2016isogeometric} & direct or Krylov\\
This work & $C^1$ quadratics; midpoints & arbitrary graded tensor grids & real, simple, nonpositive on every grid (Theorem~\ref{thm:spectrum}) & fast diagonalization on any tensor grid; DCT/DST on uniform grids\\
\bottomrule
\end{tabular}
\end{table}

%==========================================================================
\section{Midpoint quadratic-spline collocation}\label{sec:method}
%==========================================================================
\subsection{Nodal form of Haar wavelet collocation}
Let $a=t_0<t_1<\dots<t_N=b$, with $N\ge2$, be the grid, with cells $I_k=[t_{k-1},t_k]$ of width $h_k$,
midpoints $m_k$, and midpoint spacings $d_k=(h_k+h_{k+1})/2$. HWCM expands $u''$ in Haar functions and
integrates twice. On a given grid, the Haar functions span exactly the functions that are constant on each
cell, so the Haar coefficients carry the same information as the cell values $c_k=u''|_{I_k}$.
Integrating twice gives piecewise quadratics with continuous slopes, that is, $C^1$ quadratic splines.
We use the values $U_k=u(m_k)$ at the cell midpoints, which are the HWCM collocation values, as unknowns.

\begin{theorem}[Three-term identity]\label{thm:identity}
Let $s$ be a $C^1$ quadratic spline on the grid, with $U_k=s(m_k)$ and $c_k=s''|_{I_k}$. For $2\le k\le N-1$,
\begin{equation}\label{eq:interior}
\frac{U_{k+1}-U_k}{d_k}-\frac{U_k-U_{k-1}}{d_{k-1}}
=\frac{h_{k-1}^2}{8d_{k-1}}c_{k-1}+\Big(h_k-\frac{h_k^2}{8d_{k-1}}-\frac{h_k^2}{8d_k}\Big)c_k+\frac{h_{k+1}^2}{8d_k}c_{k+1}.
\end{equation}
At the left end, with $g_0=s(a)$,
\begin{align}
\frac{U_2-U_1}{d_1}-\frac{2(U_1-g_0)}{h_1}&=\Big(\frac{3h_1}{4}-\frac{h_1^2}{8d_1}\Big)c_1+\frac{h_2^2}{8d_1}c_2 &&\text{(value prescribed)},\label{eq:dirrow}\\
\frac{U_2-U_1}{d_1}&=\Big(h_1-\frac{h_1^2}{8d_1}\Big)c_1+\frac{h_2^2}{8d_1}c_2 &&\text{(if } s'(a)=0),\label{eq:neurow}
\end{align}
and symmetrically at the right end. In matrix form, $LU+b=Ac$, where $L$ is a tridiagonal difference
matrix, $A$ a tridiagonal averaging matrix, and $b$ contains the prescribed end values.
\end{theorem}

\begin{corollary}[Invertibility and solvability]\label{cor:ops}
For every grid and every combination of end conditions, $A$ has a positive diagonal and positive
off-diagonal entries, and it is strictly diagonally dominant by columns, hence invertible. With zero-slope ends, the columns of $A$
sum to the cell widths $h_k$. The collocation problem $s''(m_k)=f(m_k)$, $k=1,\dots,N$, has a
unique solution if a value is prescribed at one or both ends. If the slope vanishes at both ends, the
constants form the kernel of $L$, and a solution exists if and only if $\sum_kh_kf(m_k)=0$; it is then
unique up to an additive constant, which can be fixed, for example, by prescribing $\sum_kh_kU_k$.
\end{corollary}

The second derivative at the nodes is therefore $c=D_2U+e$, where $D_2=A^{-1}L$ and the vector $e$ collects the boundary terms.
Instead of the dense integration matrices of classical HWCM, only two tridiagonal matrices need to be
assembled. Slopes follow in closed form as well, which gives the first-derivative matrix $D_1$ and the
end slopes needed for traction conditions (Algorithm~1). Proofs of Theorem~\ref{thm:identity} and
Corollary~\ref{cor:ops} are given in Appendix~\ref{app:proofs}.

\begin{algorithm}[t]\small\hrule\smallskip
\noindent\textbf{Algorithm 1: one-dimensional operators} (\texttt{ihaar/core.py}, class \texttt{Grid1D}).
Input: knots $t_0<\dots<t_N$ and the type of each end (value prescribed, or zero slope).
\step{1} Compute $h_k$, $m_k$ and $d_k$, and assemble the tridiagonal matrices $L$ and $A$ from~\eqref{eq:interior}--\eqref{eq:neurow}.
\step{2} Factorize $A$ (tridiagonal LU) and form $D_2=A^{-1}L$ and the boundary vectors $e_i=A^{-1}b_i$.
\step{3} Form the first-derivative matrix $D_1=\mathbf 1r^{\mathsf T}+TD_2$, where $T_{kj}=h_j$ for $j<k$, $T_{kk}=h_k/2$, and $r$ encodes the end slope $\sigma_0=2(U_1-g_0)/h_1-c_1h_1/4$.
\step{4} Form the end slopes $s'(a)=\sigma_0$ and $s'(b)=\sigma_0+\sum_kh_kc_k$ as row vectors, for traction conditions.
\smallskip\hrule
\end{algorithm}

\subsection{Main result: a real, simple and nonpositive spectrum on every grid}
\begin{theorem}[Spectrum on every grid]\label{thm:spectrum}
Let $a=t_0<t_1<\dots<t_N=b$, with $N\ge2$, be any strictly increasing knot vector, and let either the
value or a zero slope be prescribed at each end. Then $D_2=A^{-1}L$ has $N$ real and simple eigenvalues,
and it is diagonalizable with a real eigenvector matrix $V$. If a value is prescribed at one or both ends,
all eigenvalues are negative. If the slope vanishes at both ends, exactly one eigenvalue is zero, with the
constant vector as eigenvector, and all others are negative.
\end{theorem}

The proof (Appendix~\ref{app:proofs}) uses only three properties of the identity: $A$ is a positive,
column-diagonally-dominant tridiagonal matrix, $-L$ is a tridiagonal M-matrix, and the columns of $A$ sum
to the cell widths. These properties make $A$ and $(-L)^{-1}$ totally nonnegative, and their product
oscillatory in the sense of Gantmacher and Krein~\cite{gantmacher2002oscillation}, whose eigenvalues are
real, positive and simple. No assumption on the grading, such as a bounded ratio between neighbouring
cells, is needed. The result extends immediately to tensor-product operators.

\begin{corollary}[Tensor-product operators]\label{cor:tensor}
On a tensor grid in $d$ dimensions, with end conditions as in Theorem~\ref{thm:spectrum} in each direction,
every operator $\sum_m\alpha_mD_2^{(m)}+\gamma I$, with $\alpha_m>0$ and $D_2^{(m)}$ acting along direction
$m$, is diagonalized by $V_1\otimes\dots\otimes V_d$, with the real eigenvalues
$\sum_m\alpha_m\lambda^{(m)}_{i_m}+\gamma$. It is nonsingular whenever $\gamma<0$, or $\gamma=0$ and a value
is prescribed at some end in at least one direction.
\end{corollary}

Theorem~\ref{thm:spectrum} is an algebraic statement: fast diagonalization is well defined on any grid,
without complex arithmetic or defective matrices. It does not bound the conditioning $\kappa(V)$ of the
eigenvector matrix, and it says nothing about the conditioning of the discrete problem itself. These
floating-point questions are examined in Section~\ref{sec:robust}.

\subsection{Uniform grids: cosine transforms, fourth order and mass conservation}
\begin{proposition}[Uniform grids]\label{thm:fourth}
On a uniform grid, $D_2=\frac{8}{h^2}\delta^2(8I+\delta^2)^{-1}$, where $\delta^2$ is the three-point second
difference with end entries $-3$ (prescribed value) or $-1$ (zero slope). It is symmetric and diagonalized
by the discrete sine (prescribed values) or cosine (zero slope) transform, with eigenvalues in
$[-8/h^2,0]$. The corrected operator
\begin{equation}\label{eq:D24}
D_2^{(4)}=D_2+\tfrac{h^2}{24}D_2^2
\end{equation}
has the same eigenvectors, is negative semidefinite, and is fourth-order accurate for smooth solutions
of zero-flux or periodic problems whose odd normal derivatives vanish at the ends. For general prescribed
values, the order remains two.
\end{proposition}

\begin{proposition}[Exact mass conservation]\label{prop:mass}
With zero slope at both ends, $h^{\mathsf T}D_2=0$ on any grid, where $h=(h_k)$, so the weighted sum
$\sum_kh_kU_k$ is conserved by evolutions of the form $\dot U=D_2W$. On a uniform grid, this sum equals the
exact integral of the spline.
\end{proposition}

Plain HWCM has half the error constant of the three-point difference ($-h^2/24$ against $h^2/12$), so it is
about twice as accurate on the same grid. The fourth-order correction costs nothing, because in the
transform domain it only changes the eigenvalues. Proofs of Propositions~\ref{thm:fourth}
and~\ref{prop:mass} are given in the Supplementary Material.

